\IfFileExists{stacks-project.cls}{%
\documentclass{stacks-project}
}{%
\documentclass{amsart}
}

% For dealing with references we use the comment environment
\usepackage{verbatim}
\newenvironment{reference}{\comment}{\endcomment}
%\newenvironment{reference}{}{}
\newenvironment{slogan}{\comment}{\endcomment}
\newenvironment{history}{\comment}{\endcomment}

% For commutative diagrams we use Xy-pic
\usepackage[all]{xy}

% We use 2cell for 2-commutative diagrams.
\xyoption{2cell}
\UseAllTwocells

% We use multicol for the list of chapters between chapters
\usepackage{multicol}

% This is generally recommended for better output
\usepackage{lmodern}
\usepackage[T1]{fontenc}

% For cross-file-references
\usepackage{xr-hyper}

% Package for hypertext links:
\usepackage{hyperref}

% For any local file, say "hello.tex" you want to link to please
% use \externaldocument[hello-]{hello}
\externaldocument[introduction-]{introduction}
\externaldocument[conventions-]{conventions}
\externaldocument[sets-]{sets}
\externaldocument[categories-]{categories}
\externaldocument[topology-]{topology}
\externaldocument[sheaves-]{sheaves}
\externaldocument[sites-]{sites}
\externaldocument[stacks-]{stacks}
\externaldocument[fields-]{fields}
\externaldocument[algebra-]{algebra}
\externaldocument[brauer-]{brauer}
\externaldocument[homology-]{homology}
\externaldocument[derived-]{derived}
\externaldocument[simplicial-]{simplicial}
\externaldocument[more-algebra-]{more-algebra}
\externaldocument[smoothing-]{smoothing}
\externaldocument[modules-]{modules}
\externaldocument[sites-modules-]{sites-modules}
\externaldocument[injectives-]{injectives}
\externaldocument[cohomology-]{cohomology}
\externaldocument[sites-cohomology-]{sites-cohomology}
\externaldocument[dga-]{dga}
\externaldocument[dpa-]{dpa}
\externaldocument[sdga-]{sdga}
\externaldocument[hypercovering-]{hypercovering}
\externaldocument[schemes-]{schemes}
\externaldocument[constructions-]{constructions}
\externaldocument[properties-]{properties}
\externaldocument[morphisms-]{morphisms}
\externaldocument[coherent-]{coherent}
\externaldocument[divisors-]{divisors}
\externaldocument[limits-]{limits}
\externaldocument[varieties-]{varieties}
\externaldocument[topologies-]{topologies}
\externaldocument[descent-]{descent}
\externaldocument[perfect-]{perfect}
\externaldocument[more-morphisms-]{more-morphisms}
\externaldocument[flat-]{flat}
\externaldocument[groupoids-]{groupoids}
\externaldocument[more-groupoids-]{more-groupoids}
\externaldocument[etale-]{etale}
\externaldocument[chow-]{chow}
\externaldocument[intersection-]{intersection}
\externaldocument[pic-]{pic}
\externaldocument[weil-]{weil}
\externaldocument[adequate-]{adequate}
\externaldocument[dualizing-]{dualizing}
\externaldocument[duality-]{duality}
\externaldocument[discriminant-]{discriminant}
\externaldocument[derham-]{derham}
\externaldocument[local-cohomology-]{local-cohomology}
\externaldocument[algebraization-]{algebraization}
\externaldocument[curves-]{curves}
\externaldocument[resolve-]{resolve}
\externaldocument[models-]{models}
\externaldocument[functors-]{functors}
\externaldocument[equiv-]{equiv}
\externaldocument[pione-]{pione}
\externaldocument[etale-cohomology-]{etale-cohomology}
\externaldocument[proetale-]{proetale}
\externaldocument[relative-cycles-]{relative-cycles}
\externaldocument[more-etale-]{more-etale}
\externaldocument[trace-]{trace}
\externaldocument[crystalline-]{crystalline}
\externaldocument[spaces-]{spaces}
\externaldocument[spaces-properties-]{spaces-properties}
\externaldocument[spaces-morphisms-]{spaces-morphisms}
\externaldocument[decent-spaces-]{decent-spaces}
\externaldocument[spaces-cohomology-]{spaces-cohomology}
\externaldocument[spaces-limits-]{spaces-limits}
\externaldocument[spaces-divisors-]{spaces-divisors}
\externaldocument[spaces-over-fields-]{spaces-over-fields}
\externaldocument[spaces-topologies-]{spaces-topologies}
\externaldocument[spaces-descent-]{spaces-descent}
\externaldocument[spaces-perfect-]{spaces-perfect}
\externaldocument[spaces-more-morphisms-]{spaces-more-morphisms}
\externaldocument[spaces-flat-]{spaces-flat}
\externaldocument[spaces-groupoids-]{spaces-groupoids}
\externaldocument[spaces-more-groupoids-]{spaces-more-groupoids}
\externaldocument[bootstrap-]{bootstrap}
\externaldocument[spaces-pushouts-]{spaces-pushouts}
\externaldocument[spaces-chow-]{spaces-chow}
\externaldocument[groupoids-quotients-]{groupoids-quotients}
\externaldocument[spaces-more-cohomology-]{spaces-more-cohomology}
\externaldocument[spaces-simplicial-]{spaces-simplicial}
\externaldocument[spaces-duality-]{spaces-duality}
\externaldocument[formal-spaces-]{formal-spaces}
\externaldocument[restricted-]{restricted}
\externaldocument[spaces-resolve-]{spaces-resolve}
\externaldocument[formal-defos-]{formal-defos}
\externaldocument[defos-]{defos}
\externaldocument[cotangent-]{cotangent}
\externaldocument[examples-defos-]{examples-defos}
\externaldocument[algebraic-]{algebraic}
\externaldocument[examples-stacks-]{examples-stacks}
\externaldocument[stacks-sheaves-]{stacks-sheaves}
\externaldocument[criteria-]{criteria}
\externaldocument[artin-]{artin}
\externaldocument[quot-]{quot}
\externaldocument[stacks-properties-]{stacks-properties}
\externaldocument[stacks-morphisms-]{stacks-morphisms}
\externaldocument[stacks-limits-]{stacks-limits}
\externaldocument[stacks-cohomology-]{stacks-cohomology}
\externaldocument[stacks-perfect-]{stacks-perfect}
\externaldocument[stacks-introduction-]{stacks-introduction}
\externaldocument[stacks-more-morphisms-]{stacks-more-morphisms}
\externaldocument[stacks-geometry-]{stacks-geometry}
\externaldocument[moduli-]{moduli}
\externaldocument[moduli-curves-]{moduli-curves}
\externaldocument[examples-]{examples}
\externaldocument[exercises-]{exercises}
\externaldocument[guide-]{guide}
\externaldocument[desirables-]{desirables}
\externaldocument[coding-]{coding}
\externaldocument[obsolete-]{obsolete}
\externaldocument[fdl-]{fdl}
\externaldocument[index-]{index}

% Theorem environments.
%
\theoremstyle{plain}
\newtheorem{theorem}[subsection]{Theorem}
\newtheorem{proposition}[subsection]{Proposition}
\newtheorem{lemma}[subsection]{Lemma}

\theoremstyle{definition}
\newtheorem{definition}[subsection]{Definition}
\newtheorem{example}[subsection]{Example}
\newtheorem{exercise}[subsection]{Exercise}
\newtheorem{situation}[subsection]{Situation}

\theoremstyle{remark}
\newtheorem{remark}[subsection]{Remark}
\newtheorem{remarks}[subsection]{Remarks}

\numberwithin{equation}{subsection}

% Macros
%
\def\lim{\mathop{\mathrm{lim}}\nolimits}
\def\colim{\mathop{\mathrm{colim}}\nolimits}
\def\Spec{\mathop{\mathrm{Spec}}}
\def\Hom{\mathop{\mathrm{Hom}}\nolimits}
\def\Ext{\mathop{\mathrm{Ext}}\nolimits}
\def\SheafHom{\mathop{\mathcal{H}\!\mathit{om}}\nolimits}
\def\SheafExt{\mathop{\mathcal{E}\!\mathit{xt}}\nolimits}
\def\Sch{\mathit{Sch}}
\def\Mor{\mathop{\mathrm{Mor}}\nolimits}
\def\Ob{\mathop{\mathrm{Ob}}\nolimits}
\def\Sh{\mathop{\mathit{Sh}}\nolimits}
\def\NL{\mathop{N\!L}\nolimits}
\def\CH{\mathop{\mathrm{CH}}\nolimits}
\def\proetale{{pro\text{-}\acute{e}tale}}
\def\etale{{\acute{e}tale}}
\def\QCoh{\mathit{QCoh}}
\def\Ker{\mathop{\mathrm{Ker}}}
\def\Im{\mathop{\mathrm{Im}}}
\def\Coker{\mathop{\mathrm{Coker}}}
\def\Coim{\mathop{\mathrm{Coim}}}

% Boxtimes
%
\DeclareMathSymbol{\boxtimes}{\mathbin}{AMSa}{"02}

%
% Macros for moduli stacks/spaces
%
\def\QCohstack{\mathcal{QC}\!\mathit{oh}}
\def\Cohstack{\mathcal{C}\!\mathit{oh}}
\def\Spacesstack{\mathcal{S}\!\mathit{paces}}
\def\Quotfunctor{\mathrm{Quot}}
\def\Hilbfunctor{\mathrm{Hilb}}
\def\Curvesstack{\mathcal{C}\!\mathit{urves}}
\def\Polarizedstack{\mathcal{P}\!\mathit{olarized}}
\def\Complexesstack{\mathcal{C}\!\mathit{omplexes}}
% \Pic is the operator that assigns to X its picard group, usage \Pic(X)
% \Picardstack_{X/B} denotes the Picard stack of X over B
% \Picardfunctor_{X/B} denotes the Picard functor of X over B
\def\Pic{\mathop{\mathrm{Pic}}\nolimits}
\def\Picardstack{\mathcal{P}\!\mathit{ic}}
\def\Picardfunctor{\mathrm{Pic}}
\def\Deformationcategory{\mathcal{D}\!\mathit{ef}}
\usepackage{cleveref}
\crefname{originallemma}{original lemma}{original lemmas}
\Crefname{originallemma}{Original lemma}{Original lemmas}
\makeatletter
\newlabel{lemma-equality-morphisms-left-localization}{{04VF}{source}{}{sourcepointer-1}{}}
\newlabel{lemma-equality-morphisms-left-localization@cref}{{[originallemma][0][]04VF}{[1][0][]source}{}{}{}}
\makeatother
\crefname{originalremark}{original remark}{original remarks}
\Crefname{originalremark}{Original remark}{Original remarks}
\makeatletter
\newlabel{remark-left-localization-morphisms-colimit}{{05Q0}{source}{}{sourcepointer-2}{}}
\newlabel{remark-left-localization-morphisms-colimit@cref}{{[originalremark][0][]05Q0}{[1][0][]source}{}{}{}}
\makeatother
\title{Commutative diagrams in the localized category come from diagrams in C after enlarging along S}
\author{The Stacks Project Authors}
\date{stacks-project:categories.tex:local-snapshot-20261012; GNU FDL 1.2 or later}
\begin{document}
\maketitle
\noindent\textit{Editorial note.} Counted span is the author's own proof environment, source lines 4703-4753, following the statement at 4680-4701.\par\medskip
\noindent\textit{Editorial note.} The footnote at 4723-4728 gives an explicit alternative path for matching b = t\textasciicircum{}\{-1\}h; it lies inside the counted span.\par\medskip
\noindent\textit{Allowed standard prerequisites (editorial).}\begin{itemize}
\item Calculus of fractions: morphisms and composition. Definition of morphisms and composition in S\textasciicircum{}\{-1\}C in this chapter (left calculus of fractions section), used at proof lines 4704-4741. Writes a and b as fractions s\textasciicircum{}\{-1\}g and t\textasciicircum{}\{-1\}h and computes equivalence classes of composites as pairs.
\item LMS2 axiom. Left multiplicative system axiom LMS2, used at 4706-4707 and 4726-4727. Produces the commutative square completing f'' over t, and the alternative replacement path in the footnote.
\item Equality criterion for fractions. Lemma lemma-equality-morphisms-left-localization of this chapter, invoked at 4747-4750. Turns equality of equivalence classes of pairs into a single post-composition d with d t in S making the arrows equal.
\item Replacement remark. Remark remark-left-localization-morphisms-colimit of this chapter, invoked at 4720-4722. Licenses the repeated replacement of (t, f'') by postcomposition with d having d t in S.
\end{itemize}\medskip
\ifdefined\hypertarget\hypertarget{sourcepointer-1}{}\fi
\noindent\textit{Source reference.} Original lemma 04VF: sources/stacks-categories/pinned-original-source.tex, source lines 4521--4521. Navigation pointer only: the counted proof cites this statement by \textbackslash{}ref as a standing input; the deduction the counted proof has to make is carried out inside the counted span, and no part of that deduction is delegated to this pointer.\par\medskip
\ifdefined\hypertarget\hypertarget{sourcepointer-2}{}\fi
\noindent\textit{Source reference.} Original remark 05Q0: sources/stacks-categories/pinned-original-source.tex, source lines 4580--4580. Navigation pointer only: the counted proof cites this statement by \textbackslash{}ref as a standing input; the deduction the counted proof has to make is carried out inside the counted span, and no part of that deduction is delegated to this pointer.\par\medskip
\begin{lemma}
\label{lemma-left-localization-diagram}
Let $\mathcal{C}$ be a category. Let $S$ be a left multiplicative
system. If $f : X \to Y$, $f' : X' \to Y'$ are two morphisms of
$\mathcal{C}$ and if
$$
\xymatrix{
Q(X) \ar[d]_{Q(f)} \ar[r]_a & Q(X') \ar[d]^{Q(f')} \\
Q(Y) \ar[r]^b & Q(Y')
}
$$
is a commutative diagram in $S^{-1}\mathcal{C}$, then there exist
a morphism $f'' : X'' \to Y''$ in $\mathcal{C}$ and a commutative
diagram
$$
\xymatrix{
X \ar[d]_f \ar[r]_g & X'' \ar[d]^{f''} & X' \ar[d]^{f'} \ar[l]^s \\
Y \ar[r]^h & Y'' & Y' \ar[l]_t
}
$$
in $\mathcal{C}$ with $s, t \in S$ and $a = s^{-1}g$, $b = t^{-1}h$.
\end{lemma}

\begin{proof}
We choose maps and objects in the following way:
First write $a = s^{-1}g$ for some $s : X' \to X''$ in $S$ and
$g : X \to X''$. By LMS2 we can find $t : Y' \to Y''$ in $S$ and
$f'' : X'' \to Y''$ such that
$$
\xymatrix{
X \ar[d]_f \ar[r]_g & X'' \ar[d]^{f''} & X' \ar[d]^{f'} \ar[l]^s \\
Y & Y'' & Y' \ar[l]_t
}
$$
commutes. Now in this diagram we are going to repeatedly change our
choice of
$$
X'' \xrightarrow{f''} Y'' \xleftarrow{t} Y'
$$
by postcomposing both $t$ and $f''$ by a morphism $d : Y'' \to Y'''$
with the property that $d \circ t \in S$. According to
Remark \ref{remark-left-localization-morphisms-colimit}
we may after such a replacement assume that there exists a morphism
$h : Y \to Y''$ such that $b = t^{-1}h$ holds\footnote{Here is a
more down-to-earth way to see this:
Write $b = q^{-1}i$ for some $q : Y' \to Z$ in $S$ and some
$i : Y \to Z$. By LMS2 we can find $r : Y'' \to Y'''$ in $S$ and
$j : Z \to Y'''$ such that $j \circ q = r \circ t$. Now, set
$d = r$ and $h = j \circ i$.}. At this point we have everything
as in the lemma except that we don't know that the left square of the
diagram commutes.
But the definition of composition in $S^{-1} \mathcal{C}$ shows that
$b \circ Q\left(f\right)$ is the equivalence class of the pair
$(h \circ f : X \to Y'', t : Y' \to Y'')$ (since $b$ is the
equivalence class of the pair $(h : Y \to Y'', t : Y' \to Y'')$,
while $Q\left(f\right)$ is the equivalence class of the pair
$(f : X \to Y, \text{id} : Y \to Y)$), while
$Q\left(f'\right) \circ a$ is the equivalence class of the pair
$(f'' \circ g : X \to Y'', t : Y' \to Y'')$ (since $a$ is the
equivalence class of the pair $(g : X \to X'', s : X' \to X'')$,
while $Q\left(f'\right)$ is the equivalence class of the pair
$(f' : X' \to Y', \text{id} : Y' \to Y')$).
Since we know that
$b \circ Q\left(f\right) = Q\left(f'\right) \circ a$, we thus
conclude that the equivalence classes of the pairs
$(h \circ f : X \to Y'', t : Y' \to Y'')$ and
$(f'' \circ g : X \to Y'', t : Y' \to Y'')$ are equal.
Hence using
Lemma \ref{lemma-equality-morphisms-left-localization}
we can find a morphism $d : Y'' \to Y'''$ such that
$d \circ t \in S$ and $d \circ h \circ f = d \circ f'' \circ g$.
Hence we make one more replacement of the kind described
above and we win.
\end{proof}

\end{document}
